\documentclass[11pt]{article}

\usepackage[T1]{fontenc}
\usepackage{lmodern}
\usepackage[margin=1in]{geometry}
\usepackage{amsmath,amssymb,bm,mathtools}
\usepackage{xcolor}
\usepackage{hyperref}

\setlength{\emergencystretch}{2em}
\allowdisplaybreaks

\hypersetup{
  colorlinks=true,
  linkcolor=blue!55!black,
  citecolor=blue!55!black,
  urlcolor=blue!55!black,
  pdftitle={Fixed-beta c-Recursion for the Sphere Ramond Channel},
  pdfauthor={Private working note}
}

\newcommand{\R}{\mathrm R}
\newcommand{\NS}{\mathrm{NS}}
\newcommand{\F}{\mathcal F}
\newcommand{\Hh}{\mathcal H}
\newcommand{\E}{\mathcal E}
\newcommand{\B}{\mathcal B}
\newcommand{\Ainf}{\mathfrak A_{\R,\infty}}
\newcommand{\Res}{\operatorname*{Res}}

\title{Fixed-\(\beta\) \(c\)-Recursion for the Sphere Ramond Channel\\[0.35em]
\large Direct Ward--Gram power counting and the regular seed}
\author{Private working note}
\date{\today}

\begin{document}

\maketitle

\begin{abstract}
We study one problem only: the sphere four-point chiral block
\(\langle V^{\NS}_4(\infty)R_3(1)R_2(z)V^{\NS}_1(0)\rangle\), factorized
through a generic Ramond module.  In the central-charge recursion all
Ramond momenta \(\beta\) are held fixed, so
\(h_\R(c,\beta)=c/24-\beta^2\) and \(G_0^2=-\beta^2\).
The Ramond Kac determinant gives the moving \(c\)-plane poles, while the
singular-vector norm and factorization of the two endpoint Ward forms
give their residues directly.  The regular block at \(c=\infty\) is a
separate problem.  We first determine its allowed powers by direct
Ward--Gram counting, following the same Gram--Schmidt and normalized
three-point-function argument as the human note, and introduce the
contracted algebra only afterwards as an equivalent organization.

The direct count keeps \(h_\R=c/24-\beta^2\) throughout.  Every nonzero
Ramond mode costs \(c^{1/2}\) in unit normalization.  Pairwise leading
contractions---the central term together with the order-\(c\) \(L_0\)
background---are neutral, an unpaired mode costs \(c^{-1/2}\), while a
Virasoro mode ending singly on the order-\(c\) Ramond weight contributes
\(c^{1/2}\).  These coherent sources make the raw block polynomially
unbounded.  Only after this classification do we identify and remove the
universal background.

The all-level sewing can be evaluated without guessing the displaced
vertices.  Decompose the Ramond module into orthogonal global
\(\mathfrak{sl}_2\) families before taking the limit.  The strict R--R
intertwiner is the identity on the nonzero oscillator Fock space, and
the Möbius transport of the pair of pants preserves the global family
generated by the Ramond ground state.  Every positive-level Ramond
quasi-primary coupling is consequently \(O(c^{-1/2})\) after unit-norm
normalization, so its sewn coefficient is \(O(c^{-1})\).  Only the
ground global family survives.  Its exact hypergeometric block and
Euler's transformation give
\[
 \lim_{c\to\infty}(1-z)^{c/12}\B_c(z)
 =
 (1-z)^{h_1+h_4+\beta_2^2+\beta_3^2}.
\]
After the explicitly defined elliptic change of normalization, it gives
the all-order regular term
\[
 S_\beta(q)=
 \left(\frac{16q}{z(q)}\right)^{\beta^2}
 (1-z(q))^{h_1+h_4+1/16}
 \theta_3(q)^{
 4(h_1+h_4-\beta_2^2-\beta_3^2)+1/4}.
\]
Direct PBW/Gram/Ward sewing through level five, performed along the
fixed-\(\beta\) family, reproduces this limit for all four chiral sign
choices.
\end{abstract}

\tableofcontents

\section{The question}

Consider the even chiral block
\begin{equation}
 \left\langle
 V^{\NS}_4(\infty)\,
 R_{\beta_3}(1)\,
 R_{\beta_2}(z)\,
 V^{\NS}_1(0)
 \right\rangle
 \label{eq:four-point-ordering}
\end{equation}
in the channel in which \(R_{\beta_2}(z)V^{\NS}_1(0)\) produces a
generic long Ramond representation of momentum \(\beta\).  The analytic
data held fixed as \(c\) varies are
\begin{equation}
 h_1,\quad h_4,\quad
 \beta_2,\quad\beta_3,\quad\beta,
 \qquad
 \eta_2,\eta_3\in\{+1,-1\}.
 \label{eq:fixed-data}
\end{equation}
The Ramond weights therefore move with \(c\):
\begin{equation}
 h_\R(c,\gamma)=\frac{c}{24}-\gamma^2,
 \qquad
 \gamma\in\{\beta,\beta_2,\beta_3\}.
 \label{eq:fixed-beta-family}
\end{equation}

This is the central point of the construction.  Since
\begin{equation}
 G_0^2=L_0-\frac{c}{24},
 \label{eq:g0-square}
\end{equation}
the internal ground space obeys \(G_0^2=-\beta^2\), independently of
\(c\).  There is no moving level-zero pole.  We work at
\(\beta\ne0\); the shortened \(\beta=0\) quotient is a separate limit.

The construction has two logically independent ingredients.  At finite
\(c\), the Ramond Kac determinant and singular-vector factorization
determine the complete pole part directly on the fixed-\(\beta\)
\(c\)-plane.  At \(c=\infty\), the regular part is first constrained by
direct power counting of the exact Ward--Gram sewing.  The contracted
Ramond algebra is introduced later as a reorganization of the surviving
background sector.  The nontrivial large-\(c\) problem is
\begin{equation}
 \boxed{\text{compute the regular \(c\to\infty\) block at fixed
 \(\beta,\beta_2,\beta_3,h_1,h_4\).}}
 \label{eq:only-problem}
\end{equation}

\section{Definition of the block}

Let \(\nu_i\) denote the NS highest-weight state of weight \(h_i\):
\begin{equation}
 L_0\nu_i=h_i\nu_i,\qquad
 L_{n>0}\nu_i=G_{r>0}\nu_i=0,
 \qquad r\in\mathbb Z+\frac12.
 \label{eq:ns-highest-state}
\end{equation}
Let \(\mathcal W_{c,\gamma}\) be a Ramond Verma module.  Its even ground
state \(w_\gamma^+\) is fixed by
\begin{equation}
 L_0w_\gamma^+=h_\R(c,\gamma)w_\gamma^+,\qquad
 L_{n>0}w_\gamma^+=G_{n>0}w_\gamma^+=0,\qquad
 (-1)^F w_\gamma^+=w_\gamma^+,\qquad
 \langle w_\gamma^+|w_\gamma^+\rangle=1.
 \label{eq:r-ground-state}
\end{equation}
The other ground state is generated by \(G_0w_\gamma^+\).  We do not
normalize it separately; its norm is proportional to
\(h_\R-c/24=-\gamma^2\).

An even internal PBW basis state at level \(n\) is denoted
\begin{equation}
 w_{\beta,I}^+
 =
 L_{-m_1}\cdots L_{-m_p}\,
 G_{-k_1}\cdots G_{-k_q}\,
 G_0^{\epsilon}\,w_\beta^+,
 \label{eq:r-pbw-state}
\end{equation}
where \(m_a,k_b\in\mathbb Z_{>0}\), the fermionic mode numbers \(k_b\)
are distinct, and
\begin{equation}
 |I|=\sum_a m_a+\sum_b k_b=n,\qquad
 q+\epsilon=0\pmod 2,\qquad \epsilon\in\{0,1\}.
 \label{eq:r-pbw-label}
\end{equation}
Thus the superscript \(+\) on \(w_{\beta,I}^+\) denotes the total even
fermion-parity sector, while \(I\) is shorthand for the complete PBW
label \((m_1,\ldots,m_p;k_1,\ldots,k_q;\epsilon)\).

In Eq.~\eqref{eq:ward-gram-block} below,
\(\nu_4\in\mathcal V_{h_4}^{\NS}\),
\(w_{\beta_3}^+\in\mathcal W_{c,\beta_3}\), and
\(w_{\beta,I}^+\in\mathcal W_{c,\beta}\).  Likewise,
\(\nu_1\in\mathcal V_{h_1}^{\NS}\) and
\(w_{\beta_2}^+\in\mathcal W_{c,\beta_2}\).
At integer level \(n\), let \(B_{\R,n}(c,\beta)\) denote the Gram matrix
of the basis \eqref{eq:r-pbw-state}.

Let the unnormalized chiral three-point forms have domains
\begin{align}
 \varrho_{NR}^{(\eta_3)}&:
 \mathcal V_{h_4}^{\NS}\times
 \mathcal W_{c,\beta_3}\times
 \mathcal W_{c,\beta}\longrightarrow\mathbb C,
 \notag\\
 \varrho_{RN}^{(\eta_2)}&:
 \mathcal W_{c,\beta}\times
 \mathcal W_{c,\beta_2}\times
 \mathcal V_{h_1}^{\NS}\longrightarrow\mathbb C.
 \label{eq:three-form-domains}
\end{align}
Here \(\eta_2,\eta_3=\pm1\) choose one of the two even Ramond chiral
structures.  Their ground-state matrix elements are the chiral
three-point coefficients
\begin{equation}
 C_L^{(\eta_3)}
 :=\varrho_{NR}^{(\eta_3)}
   (\nu_4,w_{\beta_3}^+,w_\beta^+),
 \qquad
 C_R^{(\eta_2)}
 :=\varrho_{RN}^{(\eta_2)}
   (w_\beta^+,w_{\beta_2}^+,\nu_1).
 \label{eq:three-point-coefficients}
\end{equation}
They are not set to one.  The Ward identities determine descendant
matrix elements only relative to these coefficients.  Away from their
zeros, the normalized three-forms used to define the conformal
block are
\[
 \rho_{NR}^{(\eta_3)}:=
 \frac{\varrho_{NR}^{(\eta_3)}}{C_L^{(\eta_3)}},
 \qquad
 \rho_{RN}^{(\eta_2)}:=
 \frac{\varrho_{RN}^{(\eta_2)}}{C_R^{(\eta_2)}}.
\]
Thus the normalized ground values are one by definition, while
Eq.~\eqref{eq:three-point-coefficients} retains the actual three-point
data.  At a zero of either coefficient, the normalized block is
understood by analytic continuation after this formal factorization.

The local series is
\begin{equation}
 \F_\beta^{\eta_3,\eta_2}(z)
 =
 z^{h_\R(\beta)-h_\R(\beta_2)-h_1}\,
 \B_\beta^{\eta_3,\eta_2}(c;z),
 \qquad
 \B_\beta(c;z)=\sum_{n=0}^\infty B_n(c)z^n,
 \label{eq:local-block}
\end{equation}
where the \((\eta_3,\eta_2)\) labels on \(\B_\beta\) and \(B_n\) are
suppressed on the second equality, and
\begin{equation}
 B_n(c)=
 \sum_{|I|=|J|=n}
 \rho_{NR}^{(\eta_3)}
   \bigl(\nu_4,w_{\beta_3}^+,w_{\beta,I}^+\bigr)
 \bigl(B_{\R,n}^{-1}\bigr)^{IJ}
 \rho_{RN}^{(\eta_2)}
   \bigl(w_{\beta,J}^+,w_{\beta_2}^+,\nu_1\bigr).
 \label{eq:ward-gram-block}
\end{equation}
The convention in Eq.~\eqref{eq:ward-gram-block} is the following.
\begin{enumerate}
 \item The radial ordering is
 \[
  V^{\NS}_1(0),\qquad R_{\beta_2}(z),\qquad
  R_{\beta_3}(1),\qquad V^{\NS}_4(\infty).
 \]
 After the overall power of \(z\) in Eq.~\eqref{eq:local-block} is
 extracted, each normalized three-form in
 Eq.~\eqref{eq:ward-gram-block} is evaluated at unit insertion
 position.

 \item In \(\rho_{NR}(\nu_4,w_{\beta_3}^+,w_{\beta,I}^+)\), the first
 argument is the outgoing NS bra, the second is the inserted external
 Ramond ground state, and the third is the incoming internal Ramond
 state.  In
 \(\rho_{RN}(w_{\beta,J}^+,w_{\beta_2}^+,\nu_1)\), the first argument
 is the outgoing internal Ramond state, the second is the inserted
 external Ramond ground state, and the third is the incoming NS ket.
 Hence these are the ordered \(NR\) and \(RN\) forms, not an \(RR\)
 form with an NS insertion.  Equivalently, in matrix-element notation,
 \begin{align}
  \rho_{NR}^{(\eta_3)}
   (\nu_4,w_{\beta_3}^+,w_{\beta,I}^+)
  &=\frac{1}{C_L^{(\eta_3)}}
  \langle\nu_4|
    R_{\beta_3}^{(\eta_3)}(1)
   |w_{\beta,I}^+\rangle,
  \notag\\
  \rho_{RN}^{(\eta_2)}
   (w_{\beta,J}^+,w_{\beta_2}^+,\nu_1)
  &=\frac{1}{C_R^{(\eta_2)}}
  \langle w_{\beta,J}^+|
    R_{\beta_2}^{(\eta_2)}(1)
   |\nu_1\rangle.
  \label{eq:rho-matrix-elements}
 \end{align}
 Here the numerators are the unnormalized chiral three-point matrix
 elements, and \(R_{\beta_i}^{(\eta_i)}\) denotes the Ramond
 intertwiner selected by the chiral-structure label \(\eta_i\).  The
 superscript does not introduce an additional external state.

 \item The Gram matrix and its inverse are
 \begin{equation}
  B_{\R,n}[I,J]
  =
  \left\langle w_{\beta,I}^+\middle|w_{\beta,J}^+\right\rangle,
  \qquad
  \sum_J B_{\R,n}[I,J]
  \bigl(B_{\R,n}^{-1}\bigr)^{JK}
  =\delta_I{}^K,
  \label{eq:gram-convention}
 \end{equation}
 with radial BPZ pairing
 \(L_n^\dagger=L_{-n}\) and \(G_n^\dagger=G_{-n}\).
 The indices \(I,J\) run over the complete even PBW basis
 \eqref{eq:r-pbw-state} at the same integer level \(n\), including the
 unnormalized \(G_0w_\beta^+\) ground direction when required by
 parity.  Thus \((B_{\R,n}^{-1})^{IJ}\) is precisely the internal
 Ramond sewing propagator.

 \item The labels \(\eta_3,\eta_2=\pm1\) choose the two even chiral
 Ward solutions at the left and right vertices.  They are not
 fermion-parity labels on \(w^+\), and they are not the nonchiral
 \(R^\pm\) field labels.  The corresponding unnormalized ground
 matrix elements are \(C_L^{(\eta_3)}\) and \(C_R^{(\eta_2)}\) in
 Eq.~\eqref{eq:three-point-coefficients}.

 \item Equation~\eqref{eq:ward-gram-block} uses the divided forms
 \(\rho=\varrho/C\), so it defines the normalized holomorphic
 conformal block, with \(B_0(c)=1\).  The unnormalized holomorphic
 contribution in this channel is
 \(C_L^{(\eta_3)}C_R^{(\eta_2)}\F_\beta^{\eta_3,\eta_2}(z)\).
 Hence the three-point coefficients have been factored out, not set to
 one.  Constructing a full correlator additionally requires the
 antiholomorphic blocks and their allowed left--right pairing.  Any
 \(c\)-dependence, zero, or pole of \(C_LC_R\) is dynamical CFT data
 and is not part of the universal representation-theoretic
 \(c\)-recursion.
\end{enumerate}
These conventions completely specify the block used below.

\section{The \(c\)-plane poles from the Ramond Kac determinant}

Use the ordinary central charge
\begin{equation}
 c(b)=\frac32+3(b+b^{-1})^2.
 \label{eq:central-charge}
\end{equation}
For \(r,s\ge1\), \(r+s\) odd, the Ramond degenerate momentum is
\begin{equation}
 \beta_{r,s}(b)=\frac{rb+s/b}{2\sqrt2}.
 \label{eq:kac-beta}
\end{equation}
At fixed internal \(\beta\), a \(c\)-pole is obtained by solving
\begin{equation}
 \beta=\beta_{r,s}(b_*),
 \qquad
 c_*=c(b_*).
 \label{eq:c-pole}
\end{equation}
Choose \(r>s\) and keep the two roots
\begin{equation}
 b_{r,s}^{(\alpha)}(\beta)
 =
 \frac{\sqrt2\,\beta+
 \alpha\sqrt{2\beta^2-rs}}{r},
 \qquad
 c_{r,s}^{(\alpha)}(\beta)
 =
 c\!\left(b_{r,s}^{(\alpha)}(\beta)\right),
 \qquad \alpha=\pm.
 \label{eq:c-pole-roots}
\end{equation}
This counts each generic pole once: the descriptions
\((s,r,1/b)\) and \(-\beta=\beta_{r,s}(-b)\) give the same \(c\)-values
and must not be added again.

For later use define the Jacobian and shifted Ramond momentum
\begin{align}
 J_{r,s}^{(\alpha)}(\beta)
 &:=
 \left.
 \frac{\partial\beta_{r,s}}{\partial c}
 \right|_{b=b_{r,s}^{(\alpha)}}
 =
 \left.
 \frac{r-s/b^2}
 {12\sqrt2\,(b+b^{-1})(1-b^{-2})}
 \right|_{b=b_{r,s}^{(\alpha)}},
 \notag\\
 \beta_{r,s}^{\prime(\alpha)}(\beta)
 &:=
 \left.
 \frac{(-1)^s}{2\sqrt2}
 \left(rb-\frac{s}{b}\right)
 \right|_{b=b_{r,s}^{(\alpha)}} .
 \label{eq:c-pole-jacobian-shift}
\end{align}
The second line is characterized by
\(h_\R(c_*,\beta'_{r,s})=
h_\R(c_*,\beta_{r,s})+rs/2\).

We now make the two representation-theoretic factors explicit.  With
the singular vector normalized by
\(\chi_{r,s}^+=L_{-1}^{rs/2}w_{r,s}^++\cdots\), the inverse Ramond null
norm is
\begin{equation}
 \boxed{
 A^R_{r,s}(b)
 =
 -2^{rs-\frac32}
 \prod_{\substack{
  m=1-r,\ldots,r\\
  n=1-s,\ldots,s\\
  m+n\equiv0\ ({\rm mod}\ 2)\\
  (m,n)\ne(0,0),(r,s)}}
 \left(mb+\frac{n}{b}\right)^{-1}.}
 \label{eq:ramond-inverse-null-norm}
\end{equation}
The even-sublattice condition in
Eq.~\eqref{eq:ramond-inverse-null-norm} is fixed by the defining
inverse-norm limit and is already verified by the exact level-one
\((r,s)=(2,1)\) Gram matrix.

For an NS weight \(h\), introduce a momentum branch
\begin{equation}
 Q=b+b^{-1},
 \qquad
 h=\frac{Q^2-\lambda(h;b)^2}{8}.
 \label{eq:ns-lambda}
\end{equation}
The branch must be chosen once and continued consistently in \(b\).
The implementation uses
\(\lambda(h;b)=-\sqrt{Q^2-8h}\).  Define
\begin{equation}
 p_k=1-r+2k,\qquad q_\ell=1-s+2\ell,\qquad
 \mu_{k\ell}(b)=
 \frac{p_kb+q_\ell/b}{2\sqrt2},
 \quad
 \substack{0\le k\le r-1\\0\le\ell\le s-1}.
 \label{eq:fusion-index-set}
\end{equation}
Then the N--R--R fusion polynomial appearing at either end of the
internal line is
\begin{equation}
 \boxed{
 \begin{aligned}
 P^\eta_{r,s}(\beta_R,h;b)
 ={}&
 \prod_{\substack{0\le k<r,\ 0\le\ell<s\\k+\ell\ {\rm even}}}
 \left[
  \frac{\lambda(h;b)}{2\sqrt2}
  -\eta\beta_R-\mu_{k\ell}(b)
 \right]
 \\
 &\times
 \prod_{\substack{0\le k<r,\ 0\le\ell<s\\k+\ell\ {\rm odd}}}
 \left[
  \frac{\lambda(h;b)}{2\sqrt2}
  +\eta\beta_R-\mu_{k\ell}(b)
 \right],
 \qquad \eta=\pm1.
 \end{aligned}}
 \label{eq:ramond-fusion-polynomial}
\end{equation}
These products are precisely the singular-vector factorization
coefficients of the normalized \(NR\) and \(RN\) three-forms; the
three-point coefficients \(C_LC_R\) remain outside the conformal
block.

Using Eqs.~\eqref{eq:ramond-inverse-null-norm} and
\eqref{eq:ramond-fusion-polynomial}, define the complete \(c\)-plane
residue multiplier
\begin{equation}
 \mathcal K_{r,s}^{(\alpha;\eta_3,\eta_2)}(\beta)
 :=
 -\frac{16^{rs/2}A^R_{r,s}(b_\alpha)}
 {J_{r,s}^{(\alpha)}(\beta)}
 P^{\eta_3}_{r,s}(\beta_3,h_4;b_\alpha)\,
 P^{\eta_2}_{r,s}(\beta_2,h_1;b_\alpha),
 \qquad
 b_\alpha=b_{r,s}^{(\alpha)}(\beta).
 \label{eq:c-residue-kernel}
\end{equation}
Thus \(\mathcal K\) is an explicit finite product; no additional
residue coefficient is left unspecified.
The residue of the normalized elliptic conformal block is then
\begin{equation}
 \boxed{
 \underset{c=c_{r,s}^{(\alpha)}(\beta)}
 {\Res}\,
 \Hh_\beta^{\eta_3,\eta_2}(c;q)
 =
 q^{rs/2}
 \mathcal K_{r,s}^{(\alpha;\eta_3,\eta_2)}(\beta)\,
 \Hh_{\beta_{r,s}^{\prime(\alpha)}(\beta)}^{\eta_3,\eta_2}
 \left(c_{r,s}^{(\alpha)}(\beta);q\right).}
 \label{eq:c-residue}
\end{equation}
Equivalently, for the local reduced block
\(\B_\beta(c;z)\) of Eq.~\eqref{eq:local-block},
\begin{equation}
 \boxed{
 \underset{c=c_{r,s}^{(\alpha)}(\beta)}
 {\Res}\,
 \B_\beta^{\eta_3,\eta_2}(c;z)
 =
 \left(\frac{z}{16}\right)^{rs/2}
 \mathcal K_{r,s}^{(\alpha;\eta_3,\eta_2)}(\beta)\,
 \B_{\beta_{r,s}^{\prime(\alpha)}(\beta)}^{\eta_3,\eta_2}
 \left(c_{r,s}^{(\alpha)}(\beta);z\right).}
 \label{eq:local-c-residue}
\end{equation}
The factor \(16^{-rs/2}\) removes the \(16^{rs/2}\) already included
in \(\mathcal K\).  Thus the local descendant series starts its
singular-vector tail at \(z^{rs/2}\), as it must.

Indeed,
\[
 \beta-\beta_{r,s}(c)
 =
 -J_{r,s}^{(\alpha)}(\beta)
 \bigl(c-c_{r,s}^{(\alpha)}(\beta)\bigr)
 +O\!\left((c-c_{r,s}^{(\alpha)})^2\right),
\]
which explains both the minus sign and the inverse Jacobian in
Eq.~\eqref{eq:c-residue-kernel}.  The endpoint labels
\((\eta_3,\eta_2)\) are unchanged on this canonical
 positive-\(\beta\) sheet.  Equations
 \eqref{eq:c-residue-kernel}--\eqref{eq:local-c-residue} assume a
 generic simple pole: \(J_{r,s}^{(\alpha)}\ne0\) and no collision with
 another Kac pole.  Ramified or colliding loci require a separate
 limiting prescription.

No new singular-vector calculation is required.  The same residue
formula holds for the oscillator-normalized block
\(\widehat\Hh\) defined in Eq.~\eqref{eq:oscillator-normalization}:
the factor \(\E_R(c,q)\) is holomorphic and nonzero in \(c\), and its
value at the pole cancels against the identical factor multiplying
the shifted block.

\section{Faithful fixed-\texorpdfstring{\(\beta\)}{beta} power counting
before contraction}
\label{sec:precontract-power-count}

We now determine the possible powers of \(c\) directly from the exact
Ramond Gram matrices and descendant three-point functions.  This is the
Ramond analogue of the Gram--Schmidt argument in the human note.  No
oscillator prefactor or contracted algebra is introduced in this section.

\subsection{PBW Gram bounds and explicit Gram--Schmidt orthogonalization}

We now repeat the bosonic PBW argument step by step.  Write a canonical
Ramond PBW label as
\begin{equation}
 I=(\lambda;\mu;\epsilon),
 \qquad
 \lambda=(1^{a_1}2^{a_2}\cdots),
 \qquad
 \mu=(\mu_1>\cdots>\mu_f\geq1),
 \label{eq:ramond-pbw-occupation-label}
\end{equation}
where the strict partition \(\mu\) records the fermionic modes and
\(f+\epsilon\) has the chosen total parity.  The Shapovalov form is even,
so opposite total-parity sectors are orthogonal and Gram--Schmidt can be
performed in each sector separately.  In one fixed parity block set
\begin{equation}
 v_I=
 \prod_{n\geq1}L_{-n}^{a_n}
 G_{-\mu_1}\cdots G_{-\mu_f}G_0^\epsilon w_\beta^+,
 \qquad
 \nu_R(I)=\sum_{n\geq1}a_n+f.
 \label{eq:ramond-pbw-mode-count}
\end{equation}
Thus \(\nu_R\) counts every nonzero creation mode but does not count
\(G_0\).  The latter acts only in the finite ground fibre,
\begin{equation}
 b_\epsilon
 =\langle G_0^\epsilon w_\beta^+|G_0^\epsilon w_\beta^+\rangle,
 \qquad
 b_0=1,
 \qquad
 b_1=-\beta^2.
 \label{eq:ramond-ground-fibre-diagonal}
\end{equation}

The fixed-\(\beta\) family changes which modes are large.  Acting on either
ground-fibre state, the matched one-mode contractions are
\begin{align}
 \langle L_{-n}e_a|L_{-n}e_b\rangle
 &=\left(\frac{cn^3}{12}-2n\beta^2\right)(B_0)_{ab},
 \notag\\
 \langle G_{-n}e_a|G_{-n}e_b\rangle
 &=\left(\frac{cn^2}{3}-2\beta^2\right)(B_0)_{ab},
 \qquad n\geq1.
 \label{eq:precontract-one-mode-norms}
\end{align}
The displayed coefficients include both the central term and the leading
piece of \(L_0=c/24-\beta^2\).  They are the Ramond analogues of the
matched central contractions in the bosonic proof.  In particular,
\(L_{-1}\) and \(G_{-1}\) now have order-\(c\) squared norm.  This is the
essential difference from a fixed-weight Virasoro module and from a
fixed-weight NS module, where \(L_{-1}\), and respectively
\(L_{-1},G_{-1/2}\), form an unscaled global family.  At fixed \(\beta\),
only the \(G_0\) ground fibre remains unscaled.

Let
\begin{equation}
 (B_{\R,n})_{IJ}=\langle v_I|v_J\rangle
 \label{eq:ramond-pbw-gram-entry}
\end{equation}
at a fixed level and in a fixed parity block.  Commuting the positive modes
of the bra to the right gives the precise analogues of the bosonic Gram
bounds:
\begin{equation}
 \boxed{
 \begin{aligned}
 (B_{\R,n})_{II}
 &=\kappa_I c^{\nu_R(I)}
   +O\!\left(c^{\nu_R(I)-1}\right),\\
 (B_{\R,n})_{IJ}
 &=O\!\left(c^{\nu_R(I)-1}\right),
 &&\nu_R(I)=\nu_R(J),\quad I\ne J,\\
 (B_{\R,n})_{IJ}
 &=O\!\left(c^{\min(\nu_R(I),\nu_R(J))}\right),
 &&\nu_R(I)\ne\nu_R(J),
 \end{aligned}}
 \label{eq:ramond-pbw-gram-bounds}
\end{equation}
where, for the canonical ordering in
Eq.~\eqref{eq:ramond-pbw-occupation-label},
\begin{equation}
 \kappa_I
 =b_\epsilon
 \prod_{n\geq1}
  \left[a_n!\left(\frac{n^3}{12}\right)^{a_n}\right]
 \prod_{j=1}^{f}\frac{\mu_j^2}{3}.
 \label{eq:ramond-pbw-leading-diagonal}
\end{equation}
For a generic long module \(\beta\ne0\), \(\kappa_I\ne0\).

The proof of Eq.~\eqref{eq:ramond-pbw-gram-bounds} is the same triangular
commutator proof as in the bosonic case, with one superalgebra refinement.
Every power of \(c\) uses one mode from the bra and one from the ket in a
matched \(L_n\)--\(L_{-n}\) commutator or
\(G_n\)--\(G_{-n}\) anticommutator.  A mixed \(L\)--\(G\) commutator, a
mismatched mode number, or a mode ending on \(G_0\) supplies no power of
\(c\).  Hence there can be at most
\(\min(\nu_R(I),\nu_R(J))\) large contractions.  If the two mode counts
are equal, attaining that maximum requires identical bosonic occupation
numbers, identical strict fermionic partitions, and the same ground-fibre
label.  Canonical PBW ordering then gives \(I=J\).  For \(I\ne J\), at
least one contraction is nonleading and the degree drops by one.  Fermion
reordering only supplies the canonical Koszul sign; it does not change the
power of \(c\).

Equivalently, define the symmetrically rescaled Gram matrix
\begin{equation}
 \widetilde B_{IJ}(c)
 =c^{-\nu_R(I)/2}(B_{\R,n})_{IJ}c^{-\nu_R(J)/2}.
 \label{eq:ramond-rescaled-gram}
\end{equation}
Then
\begin{equation}
 \widetilde B_{II}=\kappa_I+O(c^{-1}),
 \qquad
 \widetilde B_{IJ}=
 \begin{cases}
 O(c^{-1}),&\nu_R(I)=\nu_R(J),\ I\ne J,\\
 O\!\left(c^{-|\nu_R(I)-\nu_R(J)|/2}\right),
 &\nu_R(I)\ne\nu_R(J).
 \end{cases}
 \label{eq:ramond-rescaled-gram-limit}
\end{equation}
Thus the rescaled PBW Gram matrix has the nondegenerate diagonal limit
\(\operatorname{diag}(\kappa_I)\).  This makes the power counting visible
before any contracted-algebra interpretation.

To implement the orthogonalization, order the fixed-level PBW basis first
by increasing \(\nu_R\) and then by a fixed canonical word order.  Within
one fixed parity block, the ordinary Gram--Schmidt recursion is
\begin{equation}
 \mathcal R_I
 =v_I-
 \sum_{J\prec I}
 \frac{\langle\mathcal R_J|v_I\rangle}
      {\|\mathcal R_J\|^2}\,\mathcal R_J.
 \label{eq:ramond-gram-schmidt-recursion}
\end{equation}
Induction using Eq.~\eqref{eq:ramond-pbw-gram-bounds} gives the triangular
form
\begin{equation}
 \boxed{
 \mathcal R_I
 =v_I
 +\sum_{\substack{J\prec I\\\nu_R(J)<\nu_R(I)}}
   O(c^0)v_J
 +\sum_{\substack{J\prec I\\\nu_R(J)=\nu_R(I)}}
   O(c^{-1})v_J,}
 \label{eq:ramond-gram-schmidt-triangular}
\end{equation}
and
\begin{equation}
 \boxed{
 \|\mathcal R_I\|^2
 =\kappa_Ic^{\nu_R(I)}
  \left[1+O(c^{-1})\right].}
 \label{eq:precontract-gram-power}
\end{equation}
The order-one admixtures from lower mode count do not alter the leading
norm because their norms contain fewer powers of \(c\).  Equal-count
admixtures are themselves suppressed by \(c^{-1}\).

In this basis the inverse-Gram contraction takes exactly the spectral form
used in the bosonic argument:
\begin{equation}
 \bm\rho_{L,n}B_{\R,n}^{-1}\bm\rho_{R,n}^{\mathsf T}
 =\sum_{I:\,|I|=n}
 \frac{\rho_L(\mathcal R_I)\rho_R(\mathcal R_I)}
      {\|\mathcal R_I\|^2}.
 \label{eq:ramond-orthogonal-spectral-sewing}
\end{equation}
On a plumbing graph, the same change of basis is made independently on
each internal edge.  The only additional superalgebra datum is the Koszul
sign obtained when the odd half-edge labels are put into the chosen global
contraction order; this sign does not affect the powers of \(c\).

For later vertex counting it is useful to state the corresponding result
for unit vectors.  With
\(\widetilde v_J=c^{-\nu_R(J)/2}v_J\), one has
\begin{equation}
 \widehat{\mathcal R}_I
 =\kappa_I^{-1/2}\widetilde v_I
 +\sum_{\nu_R(J)<\nu_R(I)}
  O\!\left(c^{-[\nu_R(I)-\nu_R(J)]/2}\right)\widetilde v_J
 +\sum_{\substack{J\ne I\\\nu_R(J)=\nu_R(I)}}
  O(c^{-1})\widetilde v_J.
 \label{eq:ramond-unit-gram-schmidt-state}
\end{equation}
Therefore the triangular change of PBW basis cannot hide an enhancement
of a normalized three-point tensor.

As a minimal certificate, at level one take
\(v_1=L_{-1}w_\beta^+\) and
\(v_2=G_{-1}G_0w_\beta^+\).  Their Gram matrix is
\begin{equation}
 \begin{pmatrix}
 c/12-2\beta^2&-3\beta^2/2\\
 -3\beta^2/2&-\beta^2(c/3-2\beta^2)
 \end{pmatrix}.
 \label{eq:ramond-level-one-gram-schmidt-example}
\end{equation}
Both diagonal entries are order \(c\), whereas the off-diagonal entry is
order one.  Gram--Schmidt gives
\begin{equation}
 \mathcal R_2
 =G_{-1}G_0w_\beta^+
 +\frac{3\beta^2/2}{c/12-2\beta^2}L_{-1}w_\beta^+,
 \qquad
 \|\mathcal R_2\|^2=-\frac{\beta^2c}{3}+O(c^0),
 \label{eq:ramond-level-one-orthogonal-state}
\end{equation}
which exhibits explicitly the \(O(c^{-1})\) mixing of two distinct PBW
states with equal \(\nu_R\).

\subsection{Descendant three-point powers}

The Gram power is only half of the count.  At the two sphere endpoints,
the exact Virasoro Ward identities give
\begin{align}
 \rho_L(L_{-n}v)
 &=\left(h_\R(\beta)+|v|+n h_\R(\beta_3)-h_4\right)\rho_L(v),
 \notag\\
 \rho_R(L_{-n}v)
 &=\left(h_\R(\beta)+|v|+n h_\R(\beta_2)-h_1\right)\rho_R(v)
 \notag\\
 &=\left(\frac{n+1}{24}c+O(c^0)\right)\rho_R(v),
 \label{eq:precontract-virasoro-source}
\end{align}
and the same leading expression holds on the left.  This is the unusual
fixed-\(\beta\) effect: the internal and external R weights all move as
\(c/24\).  A Virasoro mode may therefore end singly on an order-\(c\)
background weight.  A supercurrent has no analogous one-mode source,
because its Ward identity ends on \(G_0\) or on an external ground-fibre
matrix and
\begin{equation}
 G_0^2=-\beta^2=O(c^0).
 \label{eq:precontract-finite-g0}
\end{equation}
Two counted modes can instead close through a matched commutator or
anticommutator.  Its central term, together with the contribution of
\(L_0=c/24-\beta^2\) when a zero mode is produced, supplies one power of
\(c\).  It follows immediately that
an unnormalized descendant three-point numerator obeys
\begin{equation}
 \rho_{L,R}(\mathcal R_I)
 =O\!\left(c^{p(I)}\right),
 \qquad
 p(I)\leq \ell(I)+\left\lfloor\frac{f(I)}2\right\rfloor .
 \label{eq:precontract-raw-three-point-bound}
\end{equation}
In particular, division by the Gram norm does not by itself suppress a
Ramond descendant.

To repeat the human-note counting term by term, choose one term in the
exact Ward reduction and partition its \(\nu_R\) counted modes into
\begin{equation}
 \begin{gathered}
 s=\text{Virasoro modes ending singly on an order-\(c\) R weight},\\
 2r=\text{modes exhausted in \(r\) pairwise order-\(c\) contractions},\\
 u=\text{unpaired modes ending on finite data}.
 \end{gathered}
 \label{eq:precontract-sru-definition}
\end{equation}
Then
\begin{equation}
 \nu_R=s+2r+u.
 \label{eq:precontract-sru-partition}
\end{equation}
The numerator supplies \(c^{s+r}\), while unit normalization supplies
\(c^{-\nu_R/2}\).  Hence one normalized endpoint tensor scales as
\begin{equation}
 \boxed{
 \rho_{L,R}^{(s,r,u)}(\widehat{\mathcal R}_I)
 =O\!\left(c^{s+r-\nu_R/2}\right)
 =O\!\left(c^{(s-u)/2}\right).}
 \label{eq:precontract-normalized-three-point-count}
\end{equation}
This is the desired faithful count.  A matched order-\(c\) pair is neutral
after normalization, each singly sourced Virasoro mode contributes
\(+\tfrac12\), and each unpaired mode contributes \(-\tfrac12\).  The
fixed-\(\beta\) source term is precisely what is absent from the ordinary
fixed-weight argument in the human note.

\subsection{Sewing and the regular powers}

For the sphere channel, sewing two normalized endpoint tensors gives
\begin{equation}
 c^{\Delta_e},
 \qquad
 \boxed{
 \Delta_e=\frac{s_L+s_R-u_L-u_R}{2}.}
 \label{eq:precontract-sphere-power-balance}
\end{equation}
Therefore
\begin{equation}
 \begin{array}{c|c}
 \Delta_e>0&\text{term in the raw subtraction polynomial},\\
 \Delta_e=0&\text{finite background survivor},\\
 \Delta_e<0&\text{vanishing contribution}.
 \end{array}
 \label{eq:precontract-power-classification}
\end{equation}
Pair-only patterns have \(s=u=0\) and give the Ramond oscillator analogue
of the human note's vacuum block.  Source patterns generate positive powers
of \(c\).  A locally unpaired mode is suppressed unless a source at the
other endpoint compensates it.

At a fixed coefficient, the exact global-family decomposition used below
organizes these patterns without changing their powers.  The part with no
unpaired oscillator is the universal background Gram covector.  At positive
quasi-primary level \(m\), its Gram-dual vector lies in
\begin{equation}
 L_{-1}\mathcal W^{\rm even}_{c,\beta,m-1}.
 \label{eq:precontract-background-image}
\end{equation}
This follows by simultaneous induction in level and mode number: removing
the leftmost descendant with the exact Ward identity produces either the
single source or the matched pair above, and the same order-\(c\) term is
obtained by commuting the adjoint mode through the Gram product.  After
subtracting this Gram component, at least one normalized oscillator is
unpaired.

If \(\chi_m\) is quasi-primary, BPZ invariance annihilates the background
piece exactly:
\begin{equation}
 \langle L_{-1}u_{m-1}|\chi_m\rangle
 =\langle u_{m-1}|L_1\chi_m\rangle=0.
 \label{eq:precontract-background-annihilation}
\end{equation}
Equation~\eqref{eq:precontract-normalized-three-point-count} then gives
\begin{equation}
 \rho_{L,R}(\widehat\chi_m)=O(c^{-1/2}),
 \qquad
 K_m(c)=O(c^{-1}),
 \qquad m\geq1,
 \label{eq:precontract-positive-family-suppression}
\end{equation}
for generic nonzero \(\beta\).  This is obtained before introducing the
contracted algebra.  The explicit global-family sum below then identifies
the surviving background and evaluates it.

For later higher-genus use, the identical bookkeeping can be stated on a
complete plumbing graph.  At a trinion \(v\), let \(s_v,r_v,u_v\) be the
numbers of sources, matched order-\(c\) pairs, and unpaired counted modes
across all incident internal legs.  If \(\nu_{v,e}\) is the number of nonzero modes
from leg \(e\), then
\begin{align}
 \sum_{e\ni v}\nu_{v,e}&=s_v+2r_v+u_v,
 \notag\\
 \rho_v^{\rm normalized}&=O(c^{d_v}),
 &d_v&=\frac{s_v-u_v}{2}.
 \label{eq:precontract-vertex-power}
\end{align}
For a fixed assignment of orthogonal states on a trivalent graph
\(\Gamma\), the total power is
\begin{equation}
 \boxed{
 \Delta_\Gamma
 =\sum_{v\in V(\Gamma)}d_v
 =\frac12\sum_{v\in V(\Gamma)}(s_v-u_v).}
 \label{eq:precontract-graph-power}
\end{equation}
The raw regular datum at each plumbing multidegree is the sum of all
configurations with \(\Delta_\Gamma\geq0\), taking connected contractions
before discarding locally suppressed terms.  Unlike the primary-endpoint
sphere channel, descendant-valued surface states can pair a positive-level
R quasi-primary with oscillators from adjacent edges.  Hence the sphere
suppression cannot be promoted to a universal higher-genus product without
performing this simultaneous graph-level count.

\section{The contracted Ramond algebra}

\subsection{Finite zero modes}

Define
\begin{equation}
 D=L_0-\frac{c}{24}.
 \label{eq:D}
\end{equation}
On the ground doublet,
\begin{equation}
 D|\beta,a\rangle=-\beta^2|\beta,a\rangle,
 \qquad
 G_0^2=D.
 \label{eq:ground}
\end{equation}
Both \(D\) and \(G_0\) remain unscaled.

\subsection{Normalized nonzero modes}

For \(n\ge1\), introduce
\begin{equation}
 a_n=\sqrt{\frac{12}{cn^3}}L_n,
 \qquad
 a_n^\dagger=\sqrt{\frac{12}{cn^3}}L_{-n},
 \qquad
 \psi_n=\sqrt{\frac{3}{cn^2}}G_n,
 \qquad
 \psi_n^\dagger=\sqrt{\frac{3}{cn^2}}G_{-n}.
 \label{eq:normalized-modes}
\end{equation}
On fixed-level states the \(c\to\infty\) brackets are
\begin{align}
 [a_n,a_m^\dagger]&=\delta_{nm},
 &
 \{\psi_n,\psi_m^\dagger\}&=\delta_{nm},
 \notag\\
 [a_n,a_m]&=[a_n^\dagger,a_m^\dagger]=0,
 &
 \{\psi_n,\psi_m\}
 &=\{\psi_n^\dagger,\psi_m^\dagger\}=0.
 \label{eq:free-oscillators}
\end{align}
The surviving zero-mode action is
\begin{align}
 [D,a_n]&=-na_n,
 &
 [D,a_n^\dagger]&=na_n^\dagger,
 \notag\\
 [D,\psi_n]&=-n\psi_n,
 &
 [D,\psi_n^\dagger]&=n\psi_n^\dagger,
 \notag\\
 [G_0,a_n]&=-\sqrt n\,\psi_n,
 &
 [G_0,a_n^\dagger]&=\sqrt n\,\psi_n^\dagger,
 \notag\\
 \{G_0,\psi_n\}&=\sqrt n\,a_n,
 &
 \{G_0,\psi_n^\dagger\}&=\sqrt n\,a_n^\dagger.
 \label{eq:zero-mode-action}
\end{align}
Together with \(G_0^2=D\), this defines
\begin{equation}
 \Ainf=
 \left(\mathfrak h_{\rm bos}\oplus\mathfrak h_{\rm fer}\right)
 \rtimes\langle D,G_0\rangle.
 \label{eq:contracted-algebra}
\end{equation}
It is an infinite free super-oscillator algebra with a common Ramond
supercharge.

\subsection{Action on states and field intertwiners}

The commutators above can be integrated to an explicit action on the
contracted Ramond Fock module.  Let \(\gamma_\beta\) be the restriction
of \(G_0\) to the ground doublet, so that
\(\gamma_\beta^2=-\beta^2\).  On the graded tensor product of this
doublet with the oscillator Fock space,
\begin{align}
 D&=-\beta^2+
 \sum_{n\geq1}n
 \left(a_n^\dagger a_n+\psi_n^\dagger\psi_n\right),
 \notag\\
 G_0&=\gamma_\beta+
 \sum_{n\geq1}\sqrt n\,
 \left(a_n^\dagger\psi_n+\psi_n^\dagger a_n\right).
 \label{eq:explicit-contracted-representation}
\end{align}
The graded tensor product is essential: \(\gamma_\beta\) anticommutes
with the fermionic oscillators.  Equation~\eqref{eq:explicit-contracted-representation}
then gives \(G_0^2=D\) and the brackets
Eq.~\eqref{eq:zero-mode-action}.  For example, if \(|0,g\rangle\) is
an oscillator vacuum with an open ground index,
\begin{align}
 G_0a_n^\dagger|0,g\rangle
 &=
 a_n^\dagger\gamma_\beta|0,g\rangle
 +\sqrt n\,\psi_n^\dagger|0,g\rangle,
 \notag\\
 G_0\psi_n^\dagger|0,g\rangle
 &=
 -\psi_n^\dagger\gamma_\beta|0,g\rangle
 +\sqrt n\,a_n^\dagger|0,g\rangle.
 \label{eq:g0-one-particle-action}
\end{align}
Thus \(G_0\) does not act as a scalar on an excited state: it exchanges
the bosonic and fermionic excitation of each mode number.

It is equally useful to record the action on fields.  First use the
R--R radial presentation, in which an even NS primary
\(\mathcal V_h(z)\) and its odd upper component
\(\mathcal U_h(z)=(G_{-1/2}\mathcal V_h)(z)\) are intertwiners between
two Ramond modules.  Superconformal covariance gives, for \(m\in\mathbb Z\),
\begin{align}
 L_m^{\rm out}\mathcal V_h-\mathcal V_hL_m^{\rm in}
 &=
 z^m\left(z\partial_z+(m+1)h\right)\mathcal V_h,
 \notag\\
 G_m^{\rm out}\mathcal V_h-\mathcal V_hG_m^{\rm in}
 &=
 z^{m+1/2}\mathcal U_h,
 \notag\\
 G_m^{\rm out}\mathcal U_h+\mathcal U_hG_m^{\rm in}
 &=
 z^{m-1/2}
 \left(z\partial_z+(2m+1)h\right)\mathcal V_h.
 \label{eq:rr-field-action}
\end{align}
In particular, the unscaled generators act by
\begin{align}
 D^{\rm out}\mathcal V_h-\mathcal V_hD^{\rm in}
 &=(z\partial_z+h)\mathcal V_h,
 \notag\\
 G_0^{\rm out}\mathcal V_h-\mathcal V_hG_0^{\rm in}
 &=z^{1/2}\mathcal U_h,
 \notag\\
 G_0^{\rm out}\mathcal U_h+\mathcal U_hG_0^{\rm in}
 &=z^{-1/2}(z\partial_z+h)\mathcal V_h.
 \label{eq:g0-field-action}
\end{align}
For \(n\geq1\), the normalized nonzero-mode relations are
\begin{align}
 a_n^{\rm out}\mathcal V_h-\mathcal V_ha_n^{\rm in}
 &=
 \frac{2\sqrt3}{\sqrt c\,n^{3/2}}\,
 z^n\left(z\partial_z+(n+1)h\right)\mathcal V_h,
 \notag\\
 \psi_n^{\rm out}\mathcal V_h-\mathcal V_h\psi_n^{\rm in}
 &=
 \frac{\sqrt3}{\sqrt c\,n}\,z^{n+1/2}\mathcal U_h,
 \notag\\
 \psi_n^{\rm out}\mathcal U_h+\mathcal U_h\psi_n^{\rm in}
 &=
 \frac{\sqrt3}{\sqrt c\,n}\,z^{n-1/2}
 \left(z\partial_z+(2n+1)h\right)\mathcal V_h.
 \label{eq:normalized-field-action}
\end{align}
The corresponding creation-mode formulae follow by setting \(m=-n\)
in Eq.~\eqref{eq:rr-field-action}.  Since \(h\) is fixed, every
right-hand side in Eq.~\eqref{eq:normalized-field-action} vanishes in
the strict contraction.  Hence the R--R presentation of the bottom NS
field intertwines all nonzero \(a_n,\psi_n\) oscillators diagonally,
while its \(D,G_0\) action remains nontrivial.

This statement must not be applied directly to the N--R--R and
R--R--N radial vertices in the sphere block.  There the inserted field
is Ramond, the supercurrent has a square-root branch, and there is no
ordinary single-mode commutator.  The exact action is the generalized
Ramond Ward system.  To display it compactly, set
\[
 C_p(x)=\binom{x}{p},
 \qquad
 \varepsilon(\xi,\eta)=-i(-1)^{|\xi|+|\eta|+1},
\]
and write
\(\rho=\rho_{NR}(\xi,\eta_3,\eta;z)\).  For \(n\in\mathbb Z\), the two
contour expansions are
\begin{align}
 &\sum_{p\geq0}C_p\left(n+\tfrac12\right)
 z^{n+1/2-p}\,
 \rho_{NR}(\xi,G_p\eta_3,\eta;z)
 \notag\\
 &\qquad=
 \sum_{p\geq0}C_p\left(\tfrac12\right)(-z)^p\,
 \rho_{NR}(G_{p-n-1/2}\xi,\eta_3,\eta;z)
 \notag\\
 &\qquad\quad+
 \varepsilon(\xi,\eta)
 \sum_{p\geq0}C_p\left(\tfrac12\right)(-1)^p z^{1/2-p}\,
 \rho_{NR}(\xi,\eta_3,G_{n+p}\eta;z),
 \label{eq:nr-field-ward-first}\\[2mm]
 &\sum_{p\geq0}C_p\left(\tfrac12\right)z^{1/2-p}\,
 \rho_{NR}(\xi,G_{p-n}\eta_3,\eta;z)
 \notag\\
 &\qquad=
 \sum_{p\geq0}C_p\left(\tfrac12-n\right)(-z)^p\,
 \rho_{NR}(G_{p+n-1/2}\xi,\eta_3,\eta;z)
 \notag\\
 &\qquad\quad+
 \varepsilon(\xi,\eta)
 \sum_{p\geq0}C_p\left(\tfrac12-n\right)(-1)^{n+p}
 z^{1/2-n-p}\,
 \rho_{NR}(\xi,\eta_3,G_p\eta;z).
 \label{eq:nr-field-ward-second}
\end{align}
These generalized Ramond Ward identities follow directly by deforming
the supercurrent contour around the three punctures; the R--N relation
is their BPZ/Hermitian counterpart.  On fixed-level states the apparently
infinite sums truncate.  If \(\eta_3=w_{\beta_3}^+\) is a Ramond
ground state and \(\xi=\nu_4\) is the external NS primary, the \(n=0\)
case of Eq.~\eqref{eq:nr-field-ward-second} isolates the field
zero-mode action:
\begin{align}
 z^{1/2}\rho_{NR}(\nu_4,G_0w_{\beta_3}^+,\eta;z)
 &=
 \rho_{NR}(G_{-1/2}\nu_4,w_{\beta_3}^+,\eta;z)
 \notag\\
 &\quad+
 \varepsilon(\nu_4,\eta)
 \sum_{p\geq0}C_p\left(\tfrac12\right)(-1)^p z^{1/2-p}
 \rho_{NR}(\nu_4,w_{\beta_3}^+,G_p\eta;z).
 \label{eq:nr-g0-field-action}
\end{align}
Thus the \(G_0\) action on the Ramond field, the upper NS component,
and the \(G_p\) action on the internal state form one closed system.

To obtain the contracted vertex amplitudes, substitute
\begin{equation}
 G_n=n\sqrt{\frac c3}\,\psi_n,\qquad
 G_{-n}=n\sqrt{\frac c3}\,\psi_n^\dagger,
 \qquad n\geq1,
 \label{eq:ward-to-contracted-modes}
\end{equation}
into Eqs.~\eqref{eq:nr-field-ward-first} and
\eqref{eq:nr-field-ward-second}, keep \(G_0\) unscaled, perform the
bosonic coherent displacement, and only then take \(c\to\infty\).
This gives a triangular determination of the \(\Gamma_L,\Gamma_R\)
appearing in Eq.~\eqref{eq:current-algebra-block}.

There is a complementary geometric way to organize the same
calculation.  Reordering an N--R--R three-point function into the R--R
presentation changes the local coordinate and the spin frame.  For
example, the Möbius map
\begin{equation}
 f(x)=\frac{x}{x-1}
 \label{eq:radial-reordering-map}
\end{equation}
fixes \(0\) and exchanges \(1\) with \(\infty\).  Its transport operator
\(\mathcal U_f\) is fixed by
\begin{align}
 \mathcal U_f^{-1}T(x)\mathcal U_f
 &=f'(x)^2T(f(x))+\frac c{12}\{f,x\},
 \notag\\
 \mathcal U_f^{-1}G(x)\mathcal U_f
 &=\varsigma_f\,f'(x)^{3/2}G(f(x)),
 \label{eq:coordinate-spin-transport}
\end{align}
where \(\varsigma_f=\pm1\) is the chosen spin lift; the Schwarzian
vanishes for Eq.~\eqref{eq:radial-reordering-map}.  Expanding
Eq.~\eqref{eq:coordinate-spin-transport} at \(x=0\) gives the explicit
mode transport
\begin{align}
 \mathcal U_f^{-1}L_m\mathcal U_f
 &=
 \sum_{j=m}^{\infty}(-1)^j
 \binom{1-m}{j-m}L_j,
 \notag\\
 \mathcal U_f^{-1}G_m\mathcal U_f
 &=
 \varsigma_f\sum_{j=m}^{\infty}(-1)^j
 \binom{\frac12-m}{j-m}G_j.
 \label{eq:explicit-mode-transport}
\end{align}
The sums are formal puncture expansions and are applied
coefficientwise.  For a bosonic creation mode the first sum terminates,
and Eq.~\eqref{eq:explicit-mode-transport} becomes
\begin{align}
 \mathcal U_f^{-1}a_N^\dagger\mathcal U_f
 ={}&
 \sum_{r=1}^{N}(-1)^r\binom{N+1}{N-r}
 \left(\frac rN\right)^{3/2}a_r^\dagger
 -\frac{1}{N^{3/2}}a_1
 \notag\\
 &+
 \sqrt c\,\frac{N+1}{4\sqrt3\,N^{3/2}}
 +
 \frac{2\sqrt3(N+1)}{\sqrt c\,N^{3/2}}D.
 \label{eq:explicit-bosonic-transport}
\end{align}
The third term is precisely the coherent source in
Eq.~\eqref{eq:coherent-source}.  The fermionic creation mode instead
undergoes a Bogoliubov transformation:
\begin{align}
 \mathcal U_f^{-1}\psi_N^\dagger\mathcal U_f
 =\varsigma_f\Bigg[
 &\sum_{r=1}^{N}(-1)^r
 \binom{N+\frac12}{N-r}\frac rN\,\psi_r^\dagger
 \notag\\
 &+
 \sum_{s\geq1}(-1)^s
 \binom{N+\frac12}{N+s}\frac sN\,\psi_s
 +
 \frac{\sqrt3}{\sqrt c\,N}
 \binom{N+\frac12}{N}G_0
 \Bigg].
 \label{eq:explicit-fermionic-transport}
\end{align}
Finally, the ground supercharge itself transforms as
\begin{equation}
 \mathcal U_f^{-1}G_0\mathcal U_f
 =
 \varsigma_f\left[
 G_0+\sqrt{\frac c3}
 \sum_{s\geq1}(-1)^s s\binom{\frac12}{s}\psi_s
 \right].
 \label{eq:explicit-g0-transport}
\end{equation}
Thus the bosonic transport is affine, the fermionic transport mixes
creation and annihilation modes, and the Ramond zero mode remains tied
to the fermionic annihilation modes.  These transformations, together
with the finite ground-fibre intertwiner in
Eq.~\eqref{eq:g0-field-action}, reduce the regular block to a
coherent/Bogoliubov Fock-space sewing problem.

Although
the NS field is transparent to the strict nonzero oscillators in
Eq.~\eqref{eq:normalized-field-action}, \(\mathcal U_f\) has a
nontrivial coherent large-\(c\) limit because \(L_0=c/24+D\).
Consequently the simple R--R intertwiner and the original coherent
N--R--R vertex are compatible rather than contradictory.  Evaluating
this transported current-algebra intertwiner is an equivalent route
to the regular block.

The mechanism is already visible at level one.  Since
\(f'(0)=-1\) and \(f''(0)=-2\), the descendant transforms, relative
to the primary transport, as
\begin{equation}
 \mathcal U_fL_{-1}w_\beta
 =
 \left(-L_{-1}+2h_\R(\beta)\right)\mathcal U_fw_\beta.
 \label{eq:level-one-coordinate-transport}
\end{equation}
In the R--R presentation the fixed-weight NS insertion gives
\[
 \frac{\rho_{RR}(w_{\beta_3}^+,\nu_4,L_{-1}w_\beta^+)}
 {\rho_{RR}(w_{\beta_3}^+,\nu_4,w_\beta^+)}
 =
 h_\R(\beta)+h_4-h_\R(\beta_3).
\]
After Eq.~\eqref{eq:level-one-coordinate-transport}, the coefficient
in the original N--R--R presentation is
\begin{equation}
 -\left[h_\R(\beta)+h_4-h_\R(\beta_3)\right]
 +2h_\R(\beta)
 =
 h_\R(\beta)+h_\R(\beta_3)-h_4,
 \label{eq:transport-reproduces-coherent-source}
\end{equation}
which is exactly Eq.~\eqref{eq:virasoro-ward}.  Its \(O(c)\) part
therefore comes from local-coordinate transport, not from a
nontrivial strict commutator of \(a_n\) with the fixed-weight NS field.

\subsection{The regular term is the full current-algebra block}

The contracted algebra in Eq.~\eqref{eq:contracted-algebra} is the
intermediate-state algebra of the regular term.  To say that the
regular term is its four-point block means that one must sew the two
contracted three-point vertices over the \emph{complete even Fock sector},
including every bosonic oscillator \(a_n^\dagger\), every fermionic
oscillator \(\psi_n^\dagger\), and the open Ramond ground fibre.

More explicitly, order the ground zero mode to the right and introduce
the even Fock basis
\begin{equation}
 |\bm r,\bm\epsilon;\sigma\rangle
 =
 \prod_{n\geq1}(a_n^\dagger)^{r_n}
 \prod_{n\geq1}(\psi_n^\dagger)^{\epsilon_n}
 G_0^\sigma w_\beta^+,
 \qquad
 \sigma+\sum_{n\geq1}\epsilon_n=0\pmod2,
 \label{eq:contracted-fock-basis}
\end{equation}
where only finitely many entries are nonzero,
\(r_n\in\mathbb Z_{\geq0}\), and
\(\epsilon_n,\sigma\in\{0,1\}\).  Its level is
\begin{equation}
 N(\bm r,\bm\epsilon)
 =\sum_{n\geq1}n(r_n+\epsilon_n).
 \label{eq:contracted-level}
\end{equation}
With \(\langle w_\beta^+|w_\beta^+\rangle=1\), the strict contracted Gram
matrix in this ordered basis is diagonal:
\begin{equation}
 \langle\bm r,\bm\epsilon;\sigma
 |\bm r',\bm\epsilon';\sigma'\rangle
 =
 \delta_{\bm r\bm r'}\delta_{\bm\epsilon\bm\epsilon'}
 \delta_{\sigma\sigma'}
 \left(\prod_{n\geq1}r_n!\right)(-\beta^2)^\sigma.
 \label{eq:contracted-gram}
\end{equation}

Let
\(\Gamma_L^{\eta_3}(\bm r,\bm\epsilon;\sigma)\) and
\(\Gamma_R^{\eta_2}(\bm r,\bm\epsilon;\sigma)\) denote the finite
left and right vertex amplitudes obtained from the normalized
N--R--R and R--R--N Ward identities after the coherent displacement
described below.  The current-algebra block is then, by definition,
\begin{equation}
 \boxed{
 \B_{{\rm cur},\beta}^{\eta_3,\eta_2}(z)
 =
 \sum_{\substack{\bm r,\bm\epsilon,\sigma\\
                  \sigma+\sum_n\epsilon_n\equiv0\ ({\rm mod}\ 2)}}
 z^{N(\bm r,\bm\epsilon)}
 \frac{
  \Gamma_L^{\eta_3}(\bm r,\bm\epsilon;\sigma)
  \Gamma_R^{\eta_2}(\bm r,\bm\epsilon;\sigma)}
 {\left(\prod_{n\geq1}r_n!\right)(-\beta^2)^\sigma}.}
 \label{eq:current-algebra-block}
\end{equation}
Here \(\Gamma_L(\bm0,\bm0;0)=\Gamma_R(\bm0,\bm0;0)=1\), as follows
from the normalized three-point forms.  Equation~\eqref{eq:current-algebra-block}
is written for generic \(\beta\neq0\); the \(\beta=0\) ground fibre
requires a separate limiting prescription.
Thus the summation over the \(a_n\) and \(\psi_n\) excitations is not
an optional correction to the seed: it \emph{is} the regular term.
Equation~\eqref{eq:current-algebra-block} is not a vacuum character,
because the two \(\Gamma\)'s retain the external weights, external
Ramond momenta, signs \(\eta_3,\eta_2\), and all vertex Ward data.
The contracted algebra fixes the propagator
Eq.~\eqref{eq:contracted-gram}; the contracted Ward identities fix the
two vertices.

There is one important order-of-limits issue in defining the finite
\(\Gamma\)'s.  The N--R--R vertices become coherent.  For example, the
Virasoro Ward identities give
\begin{align}
 \frac{\rho_L(L_{-n}w_\beta)}{\rho_L(w_\beta)}
 &=
 h_\R(\beta)+nh_\R(\beta_3)-h_4,
 \notag\\
 \frac{\rho_R(L_{-n}w_\beta)}{\rho_R(w_\beta)}
 &=
 h_\R(\beta)+nh_\R(\beta_2)-h_1.
 \label{eq:virasoro-ward}
\end{align}
Consequently,
\begin{equation}
 \frac{\rho_{L,R}(a_n^\dagger w_\beta)}
 {\rho_{L,R}(w_\beta)}
 =
 \sqrt c\,\frac{n+1}{4\sqrt3\,n^{3/2}}
 +O(c^{-1/2}).
 \label{eq:coherent-source}
\end{equation}
Thus the raw normalized vertex amplitudes do not possess a finite
mode-by-mode limit before the coherent piece is removed.  Subleading
mode-changing brackets can combine with the \(O(\sqrt c)\) sources and
leave finite contributions to the displaced vertices.  Consequently,
the strict diagonal Gram matrix must not be inserted before the
vertices have been displaced.  The correct operation is:
\begin{enumerate}
 \item rewrite the exact \(NR/RN\) Ward identities in the normalized
 variables;
 \item perform the coherent displacement;
 \item take the coefficientwise large-\(c\) limit to obtain the finite
 amplitudes \(\Gamma_L\) and \(\Gamma_R\);
 \item perform the complete Fock-space sewing in
 Eq.~\eqref{eq:current-algebra-block}.
\end{enumerate}
Once the two displaced vertices are known, no finite-\(c\) descendant
data remain: the regular term is computed entirely as the four-point
block of \(\Ainf\).  In the local normalization used here, its relation
to the original block is
\begin{equation}
 \boxed{
 \B_{{\rm cur},\beta}^{\eta_3,\eta_2}(z)
 =
 \lim_{\substack{c\to\infty\\
 \beta,\beta_2,\beta_3,h_1,h_4\ {\rm fixed}}}
 (1-z)^{c/12}\B_\beta(c;z)
 =\B_\infty^{\rm red}(z).}
 \label{eq:regular-is-current-block}
\end{equation}
The regular function \(S_\beta(q)\) in the elliptic recursion is this
same current-algebra block after the elliptic prefactor and oscillator
normalization have been converted from \(z\) to \(q\); it is not an
independent character-like input.

\section{All-level evaluation of the local large-\(c\) block}

Set
\begin{equation}
 A=h_1+h_4+\beta_2^2+\beta_3^2.
 \label{eq:A}
\end{equation}
The coherent oscillator sources are awkward in a PBW basis because a
nominally subleading mode-changing bracket can multiply an
\(O(\sqrt c)\) source.  The way to avoid this order-of-limits problem is
to decompose the module into orthogonal global conformal families
\emph{before} taking \(c\to\infty\).

\subsection{Exact global-family decomposition}

Let
\begin{equation}
 \mathcal Q_m^{\rm even}
 =
 \left\{\chi\in\mathcal W_{c,\beta,m}^{\rm even}:L_1\chi=0\right\}
 \label{eq:ramond-quasiprimary-space}
\end{equation}
be the even Ramond quasi-primary space at level \(m\).  At generic
\((c,\beta)\), choose a basis
\(\{\chi_{m,a}\}_{a=1}^{d_m}\) and write
\begin{equation}
 Q_m[a,b]=\langle\chi_{m,a}|\chi_{m,b}\rangle .
 \label{eq:quasiprimary-gram}
\end{equation}
The standard \(\mathfrak{sl}_2\) decomposition is orthogonal:
\begin{equation}
 \mathcal W_{c,\beta,N}^{\rm even}
 =
 \bigoplus_{m=0}^{N}
 L_{-1}^{\,N-m}\mathcal Q_m^{\rm even}.
 \label{eq:sl2-family-decomposition}
\end{equation}
Indeed, if \(m>m'\), moving \(L_1^{N-m}\) through the pairing reduces
the cross term to
\begin{equation}
 \langle\chi_{m,a}|L_{-1}^{m-m'}\chi_{m',b}\rangle=0.
 \label{eq:global-family-cross-term}
\end{equation}

Define the quasi-primary endpoint couplings
\begin{align}
 \lambda_{L,m,a}
 &=
 \rho_{NR}^{(\eta_3)}
 \left(\nu_4,w_{\beta_3}^+,\chi_{m,a}\right),
 \notag\\
 \lambda_{R,m,a}
 &=
 \rho_{RN}^{(\eta_2)}
 \left(\chi_{m,a},w_{\beta_2}^+,\nu_1\right),
 \notag\\
 K_m(c)
 &=
 \sum_{a,b=1}^{d_m}
 \lambda_{L,m,a}(Q_m^{-1})^{ab}\lambda_{R,m,b}.
 \label{eq:quasiprimary-sewing-coefficient}
\end{align}
The global Ward identities and
\[
 \langle L_{-1}^k\chi_{m,a}|L_{-1}^k\chi_{m,b}\rangle
 =
 k!\,(2h_\R(\beta)+2m)_k\,Q_m[a,b]
\]
give the exact, finite-\(c\) decomposition
\begin{equation}
 \boxed{
 \B_\beta^{\eta_3,\eta_2}(c;z)
 =
 \sum_{m=0}^{\infty}
 z^m K_m(c)\,
 \mathcal G_{h_\R(\beta)+m}(z),}
 \label{eq:exact-global-family-block}
\end{equation}
where
\begin{equation}
 \mathcal G_H(z)
 =
 {}_2F_1\left(
 H+h_\R(\beta_2)-h_1,\,
 H+h_\R(\beta_3)-h_4;\,
 2H;\,z\right).
 \label{eq:global-family-block}
\end{equation}
The level-zero quasi-primary is
\(\chi_{0,1}=w_\beta^+\).  The normalization of the three-forms gives
\begin{equation}
 K_0(c)=1.
 \label{eq:ground-family-coefficient}
\end{equation}

\subsection{Suppression of every positive-level quasi-primary}

\paragraph{Lemma.}
For fixed
\(\beta,\beta_2,\beta_3,h_1,h_4\), with \(\beta\ne0\), and for every
fixed \(m\ge1\),
\begin{equation}
 \boxed{K_m(c)=O(c^{-1}).}
 \label{eq:positive-quasiprimary-suppression}
\end{equation}

\paragraph{Proof.}
Use the Gram--Schmidt basis of
Sec.~\ref{sec:precontract-power-count}.  At fixed level, rescale each
orthogonal basis vector by the square root of its leading Gram power; the
resulting Gram matrix has a finite nondegenerate limit for generic
\(\beta\ne0\).  No contracted commutator is needed for this statement: it
is the finite-dimensional Gaussian elimination implied by
Eq.~\eqref{eq:precontract-gram-power}.

The normalized three-point count
Eq.~\eqref{eq:precontract-normalized-three-point-count} shows that all
terms without an unpaired oscillator form the universal background
covector.  By Eq.~\eqref{eq:precontract-background-image}, its Gram-dual
lies in \(L_{-1}\mathcal W_{c,\beta,m-1}^{\rm even}\).  Hence, for every
\(\chi_{m,a}\in\ker L_1\),
\[
 \langle L_{-1}u|\chi_{m,a}\rangle
 =\langle u|L_1\chi_{m,a}\rangle=0.
\]
The quasi-primary projection therefore removes every source/pair pattern
with \(u=0\).  Every remaining normalized endpoint term contains at least
one unpaired nonzero mode, and so
\begin{equation}
 \rho_{L,R}\!\left(\widehat\chi_{m,a}\right)
 =O(c^{-1/2}),
 \qquad m\ge1,
 \label{eq:unit-quasiprimary-coupling}
\end{equation}
directly from Eq.~\eqref{eq:precontract-normalized-three-point-count}.
Sewing the two endpoints with the finite limiting inverse quasi-primary
Gram matrix gives \(K_m(c)=O(c^{-1})\).  This proves
Eq.~\eqref{eq:positive-quasiprimary-suppression} before any use of the
contracted algebra. \(\square\)

The use of the complete even quasi-primary space is important here.
The \(G_0\)-ground direction and the fermionic descendants participate
in the projection which removes the coherent piece; projecting the
ground fibre before forming \(K_m\) would invalidate the lemma.

\subsection{The surviving ground family}

Euler's hypergeometric transformation applied to
Eq.~\eqref{eq:global-family-block}, with
\(H=h_\R(\beta)+m\), gives
\begin{align}
 \mathcal G_{h_\R(\beta)+m}(z)
 ={}&
 (1-z)^{h_1+h_4-h_\R(\beta_2)-h_\R(\beta_3)}
 \notag\\
 &\times
 {}_2F_1\left(
 m+h_1+\beta_2^2-\beta^2,\,
 m+h_4+\beta_3^2-\beta^2;\,
 \frac c{12}-2\beta^2+2m;\,z\right)
 \notag\\
 ={}&
 (1-z)^{-c/12+A}\left[1+O(c^{-1})\right],
 \label{eq:global-family-large-c}
\end{align}
coefficientwise in \(z\).  Notice that the exponent is independent of
both the internal momentum and the quasi-primary level \(m\).

At a fixed power of \(z\), only finitely many terms of
Eq.~\eqref{eq:exact-global-family-block} occur.  Combining
Eqs.~\eqref{eq:ground-family-coefficient},
\eqref{eq:positive-quasiprimary-suppression}, and
\eqref{eq:global-family-large-c} therefore proves
\begin{equation}
 \boxed{
 \B_\beta^{\eta_3,\eta_2}(c;z)
 =
 (1-z)^{-c/12+A}
 \left[1+O(c^{-1})\right],}
 \label{eq:local-asymptotic}
\end{equation}
and hence
\begin{equation}
 \boxed{
 \B_\infty^{\rm red}(z)
 =
 \lim_{c\to\infty}(1-z)^{c/12}\B_\beta(c;z)
 =(1-z)^A.}
 \label{eq:local-seed}
\end{equation}
All limits are coefficientwise.  The error statement in
Eq.~\eqref{eq:local-asymptotic} is not an absolute estimate on the
unstripped coefficients, which grow polynomially in \(c\).

Equation~\eqref{eq:local-seed}, now a consequence rather than an
assumption, is equivalently
\begin{equation}
 \left[(1-z)\frac{d}{dz}+A\right]\B_\infty^{\rm red}(z)=0,
 \qquad
 \B_\infty^{\rm red}(0)=1.
 \label{eq:contracted-ward-equation}
\end{equation}
Writing
\(\B_\infty^{\rm red}(z)=\sum_{N\ge0}C_Nz^N\), the derived recurrence is
\begin{equation}
 N C_N+(A-N+1)C_{N-1}=0,
 \qquad
 C_N=(-1)^N\binom{A}{N}.
 \label{eq:derived-level-recurrence}
\end{equation}

\subsection{Level-one certificate}

In the even basis
\begin{equation}
 \left(L_{-1}w^+,\ G_{-1}G_0w^+\right),
\end{equation}
the Gram matrix is
\begin{equation}
 B_{\R,1}=
 \begin{pmatrix}
 2h_\R(\beta)&-\frac32\beta^2\\
 -\frac32\beta^2&
 -\beta^2\left(2h_\R(\beta)+\frac c4\right)
 \end{pmatrix}.
 \label{eq:level-one-gram}
\end{equation}
The two Ward vectors are
\begin{align}
 \bm\rho_L
 &=
 \left(
 h_\R(\beta)+h_\R(\beta_3)-h_4,\,
 -\frac{\beta^2}{2}-\eta_3\beta\beta_3
 \right),
 \notag\\
 \bm\rho_R
 &=
 \left(
 h_\R(\beta)+h_\R(\beta_2)-h_1,\,
 -\frac{\beta^2}{2}-\eta_2\beta\beta_2
 \right).
 \label{eq:level-one-ward}
\end{align}
Their full contraction gives
\begin{equation}
 B_1(c)=\bm\rho_LB_{\R,1}^{-1}\bm\rho_R^{\mathsf T}
 =
 \frac c{12}-A+O(c^{-1}),
 \label{eq:level-one-result}
\end{equation}
which is exactly the first coefficient of
\((1-z)^{-c/12+A}\).  The component signs drop out only after the
complete ground-fibre contraction.

\subsection{The short \(\beta=0\) quotient}

The proof above was formulated for a generic long module so that the
two-dimensional ground metric is invertible.  At \(\beta=0\) one must
not substitute directly into Eq.~\eqref{eq:current-algebra-block}.
Instead take the irreducible short quotient
\begin{equation}
 G_0w_0^+=0,
 \qquad
 h_\R(0)=\frac c{24},
 \label{eq:short-ground}
\end{equation}
and form its quotient Gram matrices.

The global-family proof simplifies in this quotient.  The even ground
quasi-primary space is one-dimensional, generated by \(w_0^+\), and
all basis vectors with a rightmost \(G_0w_0^+\) are absent.  The
positive-level quotient still contracts to the ordinary boson--fermion
Fock metric and has no positive-level vector annihilated by every
\(a_n,\psi_n\).  Hence the suppression lemma remains valid:
\begin{equation}
 K_m^{\rm short}(c)=O(c^{-1}),
 \qquad m\ge1.
 \label{eq:short-quasiprimary-suppression}
\end{equation}
The surviving ground global family is
\begin{align}
 \mathcal G_{c/24}^{\rm short}(z)
 ={}&
 {}_2F_1\left(
 \frac c{12}-\beta_2^2-h_1,\,
 \frac c{12}-\beta_3^2-h_4;\,
 \frac c{12};\,z\right)
 \notag\\
 ={}&
 (1-z)^{-c/12+A}
 {}_2F_1\left(
 h_1+\beta_2^2,\,
 h_4+\beta_3^2;\,
 \frac c{12};\,z\right).
 \label{eq:short-global-family}
\end{align}
Therefore
\begin{equation}
 \boxed{
 \lim_{c\to\infty}(1-z)^{c/12}
 \B_{0}^{\rm short}(c;z)
 =(1-z)^A.}
 \label{eq:short-local-seed}
\end{equation}
This is an independent quotient argument.  It does not assert equality
between the finite-\(c\) long block analytically continued to
\(\beta=0\) and the finite-\(c\) short block; subleading terms must be
computed with the quotient Gram matrices.  The two arguments establish
only that their normalized leading large-\(c\) term is the same.

\section{Coefficientwise finite-pole-free part at
\texorpdfstring{\(c=\infty\)}{c=infinity}}
\label{sec:coefficientwise-infinity}

The phrase ``regular part at infinity'' requires care in the fixed-\(\beta\)
problem.  At every fixed local level,
\begin{equation}
 F_N(c)=
 \bm\rho_{L,N}B_{\R,N}^{-1}\bm\rho_{R,N}^{\mathsf T}
 \label{eq:raw-local-level-coefficient}
\end{equation}
is a rational function of \(c\).  Its unique polynomial division is
\begin{equation}
 F_N(c)=P_N(c)+R_N(c),
 \qquad
 R_N(c)=O(c^{-1}).
 \label{eq:raw-polynomial-division}
\end{equation}
Here \(P_N\) is free of finite \(c\)-plane poles, but a nonconstant
\(P_N\) has a pole at \(c=\infty\); it is not holomorphic there.  The
terminology ``finite-pole-free part'' avoids this ambiguity.

Subtracting only the leading power of \(P_N\) is insufficient.  Moreover,
the lower-level proper remainders \(R_j=O(c^{-1})\) can be promoted to
nonnegative powers when they are multiplied by the order-\(c\) coefficients
of the coherent background.  This can be seen exactly at the first two
levels.  Put
\begin{equation}
 t=\frac c{12},
 \qquad
 A=h_1+h_4+\beta_2^2+\beta_3^2.
 \label{eq:infinity-ledger-tA}
\end{equation}
Direct polynomial division of the complete Ward--Gram coefficients gives
\begin{align}
 P_1(c)&=t-A,
 \notag\\
 P_2(c)&=\frac{(t-A)(t-A+1)}2+D_2,
 \label{eq:raw-polynomial-parts-level-two}
\end{align}
where
\begin{align}
 D_2={}&
 \beta^4-\beta^2\left(A+\frac14\right)
 +\frac{\beta}{4}
  \left(\beta_2\eta_2+\beta_3\eta_3\right)
 \notag\\
 &+\beta_2^2(\beta_3^2+h_4)
 +\beta_3^2h_1+h_1h_4
 -\frac14\beta_2\beta_3\eta_2\eta_3.
 \label{eq:raw-level-two-extra-polynomial}
\end{align}
The extra constant \(D_2\) is not contained in the coefficient of the
leading resummation \((1-z)^{-t+A}\).  If
\(F_1=P_1+R_1\), exact division also gives
\begin{equation}
 \lim_{c\to\infty}tR_1(c)=D_2.
 \label{eq:promoted-level-one-remainder}
\end{equation}
Thus an additive subtraction based only on the leading resummed expression
misses a finite term already at level two.

The correct operation is the exact multiplicative stripping
\begin{equation}
 \mathcal C(c;z)
 =(1-z)^{c/12}\B(c;z)
 =\sum_{N\ge0}C_N(c)z^N,
 \label{eq:exact-multiplicative-local-stripping}
\end{equation}
where
\begin{equation}
 C_N(c)=
 \sum_{k=0}^{N}(-1)^k\binom{c/12}{k}F_{N-k}(c),
 \qquad F_0=1.
 \label{eq:exact-stripped-coefficients}
\end{equation}
For the exact level-one and level-two Ward--Gram data,
\begin{align}
 \operatorname{Pol}_{\infty}C_1(c)&=-A,
 \notag\\
 \operatorname{Pol}_{\infty}C_2(c)&=\frac{A(A-1)}2,
 \notag\\
 C_1(c)+A&=O(c^{-1}),
 \qquad
 C_2(c)-\frac{A(A-1)}2=O(c^{-1}).
 \label{eq:stripped-polynomial-parts-level-two}
\end{align}
Equation~\eqref{eq:promoted-level-one-remainder} is included automatically
in Eq.~\eqref{eq:exact-stripped-coefficients}; it cancels the apparent
\(D_2\) in the stripped polynomial part.  This is the coefficientwise
criterion needed for the \(c\)-recursion:
\begin{equation}
 \boxed{
 C_N(c)=S_N+O(c^{-1})
 \quad\hbox{at every fixed level}.}
 \label{eq:proper-recursion-normalization-criterion}
\end{equation}
Only after this criterion is established may the constant \(S_N\) be used
as the regular seed and the proper remainder be reconstructed from its
finite Kac poles.

There is no branch point at \(c=\infty\) in this coefficientwise statement:
every \(F_N(c)\) and \(C_N(c)\) is rational in \(c\).  The all-level
resummation
\begin{equation}
 (1-z)^{-c/12}
 =\exp\!\left[-\frac c{12}\log(1-z)\right]
 \label{eq:resummed-essential-infinity}
\end{equation}
is entire at every finite \(c\) and has an essential singularity at
\(c=\infty\) for fixed generic \(z\).  The choice of \(\log(1-z)\) is a
branch choice in the \(z\)-plane, not in the \(c\)-plane.  The unbounded
sequence of coefficientwise pole orders is precisely what resums into this
essential singularity.

The symbolic ledger and its exact assertions are implemented in
\path{Code/analyze_ramond_sphere_c_infinity.py}.

\section{Conversion to the elliptic block}

Let
\begin{equation}
 q=\exp\left[-\pi\frac{K(1-z)}{K(z)}\right],
 \qquad
 \Delta_0=\frac{c-\frac32}{24}.
 \label{eq:nome}
\end{equation}
Define the elliptic block by the change of normalization
\begin{align}
 \F_\beta(z)
 ={}&
 (16q)^{h_\R(\beta)-\Delta_0-\frac1{16}}
 z^{\Delta_0-h_1-h_\R(\beta_2)+\frac1{16}}
 (1-z)^{\Delta_0-h_\R(\beta_2)-h_\R(\beta_3)}
 \notag\\
 &\times
 \theta_3(q)^{
 \frac{c-\frac32}{2}
 -4\left(h_1+h_\R(\beta_2)+h_\R(\beta_3)+h_4\right)
 +\frac12}
 \Hh_\beta(c;q).
 \label{eq:elliptic-prefactor}
\end{align}
At fixed \(\beta\), the raw \(\Hh_\beta\) is not bounded as
\(c\to\infty\).  Define
\begin{equation}
 P(q)=\prod_{n=1}^{\infty}\frac{1+q^n}{1-q^n},
 \qquad
 \E_R(c,q)=P(q)^{c/6},
 \qquad
 \widehat\Hh_\beta(c;q)=\E_R(c,q)^{-1}\Hh_\beta(c;q).
 \label{eq:oscillator-normalization}
\end{equation}
The product identity
\begin{equation}
 \theta_3(q)P(q)=(1-z(q))^{-1/4}
 \label{eq:product-identity}
\end{equation}
shows that the \(O(c)\) part of
Eqs.~\eqref{eq:elliptic-prefactor} and
\eqref{eq:oscillator-normalization} is exactly the local factor
\((1-z)^{-c/12}\).

Combining the proved local limit \eqref{eq:local-seed} with
Eqs.~\eqref{eq:elliptic-prefactor}--\eqref{eq:product-identity} gives the
normalized all-order regular term:
\begin{equation}
 \boxed{
 S_\beta(q)=
 \left(\frac{16q}{z(q)}\right)^{\beta^2}
 (1-z(q))^{h_1+h_4+\frac1{16}}
 \theta_3(q)^{
 4(h_1+h_4-\beta_2^2-\beta_3^2)+\frac14}.}
 \label{eq:closed-seed}
\end{equation}
The limiting identity is therefore
\begin{equation}
 S_\beta(q)=
 \lim_{\substack{c\to\infty\\
 \beta,\beta_2,\beta_3,h_1,h_4\ {\rm fixed}}}
 \widehat\Hh_\beta(c;q).
 \label{eq:seed-limit}
\end{equation}
The result is independent of \((\eta_3,\eta_2)\).  Its dependence on
the internal \(\beta\) comes entirely from the local-coordinate factor
\((16q/z)^{\beta^2}\).

\subsection{First two coefficients}

Define
\begin{equation}
 X=\beta^2-\beta_2^2-\beta_3^2-h_1-h_4.
\end{equation}
Expansion of Eq.~\eqref{eq:closed-seed} gives
\begin{align}
 S_\beta(q)
 &=
 1+S_1q+S_2q^2+O(q^3),
 \notag\\
 S_1&=8X-\frac12,
 \notag\\
 S_2&=
 32X^2-16\beta^2
 +12\beta_2^2+12\beta_3^2
 -4h_1-4h_4-\frac38.
 \label{eq:first-two}
\end{align}
These agree with the independent level-one and level-two
Ward--Gram calculations.

\section{Complete fixed-\(\beta\) \(c\)-recursion}

At each fixed power of \(q\), subtracting the coefficientwise
large-\(c\) regular part leaves a rational function of \(c\) that
vanishes at infinity.  Its partial-fraction expansion is therefore
fixed by Eq.~\eqref{eq:c-residue}.  As an identity of formal
\(q\)-series, the block recursion is
\begin{equation}
 \boxed{
 \begin{aligned}
 \widehat\Hh_\beta^{\eta_3,\eta_2}(c;q)
 ={}&S_\beta(q)\\
 &+
 \sum_{\substack{r>s\ge1\\r+s\ {\rm odd}}}
 \sum_{\alpha=\pm}
 \frac{
 q^{rs/2}
 \mathcal K_{r,s}^{(\alpha;\eta_3,\eta_2)}(\beta)}
 {c-c_{r,s}^{(\alpha)}(\beta)}
 \widehat\Hh_{\beta_{r,s}^{\prime(\alpha)}(\beta)}
 ^{\eta_3,\eta_2}
 \left(c_{r,s}^{(\alpha)}(\beta);q\right).
 \end{aligned}}
 \label{eq:functional-c-recursion}
\end{equation}
Here the external data
\((h_1,h_4,\beta_2,\beta_3)\) are the same in every block on the
right-hand side; only the internal momentum is shifted, and the
shifted block is evaluated at the pole value of \(c\).  Equation
\eqref{eq:functional-c-recursion} is exact, with the regular term
\(S_\beta(q)\) given by Eq.~\eqref{eq:closed-seed}.

Write
\begin{equation}
 \widehat\Hh_\beta(c;q)=
 \sum_{n=0}^{\infty}\widehat H_n(c;\beta)q^n,
 \qquad
 S_\beta(q)=\sum_{n=0}^{\infty}S_n(\beta)q^n.
\end{equation}
At order \(q^n\), only labels with \(rs/2\le n\) occur.  On the
canonical positive-\(\beta\) sheet the exact coefficient recursion is
\begin{align}
 \widehat H_n(c;\beta)
 ={}&
 S_n(\beta)
 \notag\\
 &+
 \sum_{\substack{r>s\ge1,\ r+s\ {\rm odd}\\rs/2\le n}}
 \sum_{\alpha=\pm}
 \frac{\mathcal K_{r,s}^{(\alpha;\eta_3,\eta_2)}(\beta)}
 {c-c_{r,s}^{(\alpha)}(\beta)}
 \widehat H_{n-rs/2}
 \left(
 c_{r,s}^{(\alpha)}(\beta);
 \beta_{r,s}^{\prime(\alpha)}(\beta)
 \right),
 \label{eq:closed-c-recursion}
\end{align}
where the pole position, shifted momentum, and complete residue
multiplier are given in Eqs.~\eqref{eq:c-pole-roots},
\eqref{eq:c-pole-jacobian-shift}, and
\eqref{eq:c-residue-kernel}.  The endpoint component labels are
unchanged on this canonical sheet.

\paragraph{Completeness of the pole part.}
For generic fixed \(\beta\ne0\), the pole sum in
Eqs.~\eqref{eq:functional-c-recursion} and
\eqref{eq:closed-c-recursion} is complete, coefficient by coefficient
in \(q\).  Indeed:
\begin{enumerate}
 \item the Ramond Kac determinant exhausts the zeros of the internal
 Gram determinant, namely
 \(\beta=\pm\beta_{r,s}(c)\), with \(r+s\) odd and \(rs/2\le n\);
 \item restricting to \(r>s\) and keeping both
 \(b_{r,s}^{(\alpha)}\), \(\alpha=\pm\), represents every resulting
 \(c\)-pole once;
 \item the fixed nonzero \(\beta\) ground fibre has
 \(G_0^2=-\beta^2\), so it creates no additional moving level-zero
 pole;
 \item \(C_LC_R\) has been factored out, and the oscillator factor
 \(\E_R(c,q)\) is entire and nonzero in \(c\), so neither alters the
 conformal-block pole set.
\end{enumerate}
Fusion-polynomial zeros can cancel individual residues but cannot
produce new poles.  At coincident or ramified Kac loci the displayed
simple-pole sum is understood by analytic continuation from generic
data.

Together with the all-order regular term
Eq.~\eqref{eq:closed-seed}, this completeness statement closes the
fixed-\(\beta\) \(c\)-recursion coefficient by coefficient.

\section{Checks and present status}

\subsection{Analytic direct truncation through level four}

The direct check uses the complete even PBW spaces, whose dimensions at
levels zero through four are
\[
 1,\quad2,\quad4,\quad8,\quad14.
\]
For reference, ordered bases may be chosen as
\begin{align*}
\mathcal H_0^{\rm even}:{}&
 \{w\},\\
\mathcal H_1^{\rm even}:{}&
 \{L_{-1}w,\ G_{-1}G_0w\},\\
\mathcal H_2^{\rm even}:{}&
 \{L_{-1}^2w,\ L_{-2}w,\ G_{-2}G_0w,\
    L_{-1}G_{-1}G_0w\},\\
\mathcal H_3^{\rm even}:{}&
 \{L_{-1}^3w,\ L_{-2}L_{-1}w,\ L_{-3}w,\ G_{-3}G_0w,\
 L_{-1}G_{-2}G_0w,\ L_{-1}^2G_{-1}G_0w,\\
&\hspace{34mm}L_{-2}G_{-1}G_0w,\ G_{-2}G_{-1}w\},\\
\mathcal H_4^{\rm even}:{}&
 \{L_{-1}^4w,\ L_{-2}L_{-1}^2w,\ L_{-2}^2w,\
 L_{-3}L_{-1}w,\ L_{-4}w,\ G_{-4}G_0w,\\
&\quad L_{-1}G_{-3}G_0w,\ L_{-1}^2G_{-2}G_0w,\
 L_{-1}^3G_{-1}G_0w,\ L_{-2}G_{-2}G_0w,\\
&\quad L_{-2}L_{-1}G_{-1}G_0w,\ L_{-3}G_{-1}G_0w,\
 G_{-3}G_{-1}w,\ L_{-1}G_{-2}G_{-1}w\}.
\end{align*}
Let \(B_{\R,N}\) be the Gram matrix and let
\(\bm\rho_{L,N},\bm\rho_{R,N}\) be the two Ward vectors in these
bases.  The exact level-four truncation is the analytic expression
\begin{equation}
 \B_\beta^{[4]}(c;z)
 =1+\sum_{N=1}^{4}z^N
 \bm\rho_{L,N}B_{\R,N}^{-1}\bm\rho_{R,N}^{\mathsf T}
 +O(z^5).
 \label{eq:analytic-level-four-truncation}
\end{equation}
This involves only finite matrices of sizes \(2,4,8,14\); no pole
recursion enters.

It is important that Eq.~\eqref{eq:analytic-level-four-truncation} is
the finite-\(c\) superconformal block itself, before any large-\(c\)
normalization.  If
\[
 F_N(c)=
 \bm\rho_{L,N}B_{\R,N}^{-1}\bm\rho_{R,N}^{\mathsf T},
\]
then an entirely scalar analytic form of the same coefficient is the
bordered-determinant formula
\begin{equation}
 F_N(c)=-
 \frac{\det\!\begin{pmatrix}
  0&\bm\rho_{L,N}\\
  \bm\rho_{R,N}^{\mathsf T}&B_{\R,N}
 \end{pmatrix}}
 {\det B_{\R,N}}.
 \label{eq:finite-c-bordered-coefficient}
\end{equation}
Thus the full four-point block, including its OPE power, is
\begin{equation}
 \boxed{
 \mathcal F_{\beta}^{\eta_3,\eta_2,[4]}(c;z)
 =z^{-\beta^2+\beta_2^2-h_1}
 \left(1+F_1(c)z+F_2(c)z^2+F_3(c)z^3+F_4(c)z^4
 +O(z^5)\right).}
 \label{eq:finite-c-full-level-four-block}
\end{equation}
All entries of the four Gram matrices and the eight Ward vectors are
polynomials in
\[
 c,\ \beta,\ \beta_2,\ \beta_3,\ h_1,\ h_4,\ \eta_2,\ \eta_3.
\]
They are recorded in the accompanying exact symbolic data file
\path{ramond_sphere_level4_exact_expansion.md}.  This is a
machine-readable form of Eq.~\eqref{eq:finite-c-full-level-four-block},
not a numerical implementation.  For example, at level one,
\begin{align*}
B_{\R,1}&=
\begin{pmatrix}
 \frac c{12}-2\beta^2&-\frac32\beta^2\\
 -\frac32\beta^2&2\beta^4-\frac c3\beta^2
\end{pmatrix},\\
\bm\rho_{L,1}&=
\left(\frac c{12}-\beta^2-\beta_3^2-h_4,
-\frac12\beta^2-\eta_3\beta\beta_3\right),\\
\bm\rho_{R,1}&=
\left(\frac c{12}-\beta^2-\beta_2^2-h_1,
-\frac12\beta^2-\eta_2\beta\beta_2\right).
\end{align*}
Substitution in Eq.~\eqref{eq:finite-c-bordered-coefficient} gives the
complete finite-\(c\) coefficient \(F_1(c)\); the higher-level entries
are obtained in exactly the same way from the displayed PBW bases.

Only after the finite-\(c\) block has been constructed in this way do
we compare it with the contracted block.  Define the coefficients
coefficients
\begin{equation}
 C_N(c)=
 \sum_{k=0}^{N}(-1)^k\binom{c/12}{k}
 \bm\rho_{L,N-k}B_{\R,N-k}^{-1}
 \bm\rho_{R,N-k}^{\mathsf T}.
 \label{eq:analytic-stripped-level-coefficient}
\end{equation}
Here the level-zero contraction is one.  Directly evaluating the finite
Ward vectors and Gram matrices and then taking \(c\to\infty\), with all
data in Eq.~\eqref{eq:fixed-data} held fixed, gives
\begin{align}
 \lim_{c\to\infty}C_0(c)&=1,\notag\\
 \lim_{c\to\infty}C_1(c)&=-A,\notag\\
 \lim_{c\to\infty}C_2(c)&=\frac{A(A-1)}{2},\notag\\
 \lim_{c\to\infty}C_3(c)&=-\frac{A(A-1)(A-2)}{6},\notag\\
 \lim_{c\to\infty}C_4(c)&=\frac{A(A-1)(A-2)(A-3)}{24}.
 \label{eq:analytic-level-four-result}
\end{align}
All dependence on \(\beta,\eta_2,\eta_3\) cancels in these five direct
contractions.  Thus the analytic level-four descendant calculation is
\begin{equation}
 \boxed{
 \begin{aligned}
 \lim_{c\to\infty}(1-z)^{c/12}\B_\beta^{[4]}(c;z)
 ={}&1-Az+\frac{A(A-1)}2z^2
 -\frac{A(A-1)(A-2)}6z^3\\
 &+\frac{A(A-1)(A-2)(A-3)}{24}z^4+O(z^5).
 \end{aligned}}
 \label{eq:analytic-level-four-block}
\end{equation}
This is precisely the truncation of \((1-z)^A\), but it has been
obtained here from the finite super-Virasoro Ward--Gram sewing.

\subsection{Other checks}

The fixed-\(\beta\) recursion and the independent direct descendant
implementation perform the following checks:
\begin{enumerate}
 \item the two level-one \(c\)-poles agree with the zeros of the exact
 \(2\times2\) Ramond Gram determinant;
 \item the residue multipliers agree with direct Laurent extraction from
 the Ward--Gram block at the simple Kac poles;
 \item Eq.~\eqref{eq:closed-seed} reproduces the independently derived
 \(S_1\) and \(S_2\);
 \item direct PBW/Gram/Ward sewing along a fixed-\(\beta\) \(c\)-family
 agrees with the closed \(c\)-recursion through level five, for all four
 choices of \((\eta_3,\eta_2)\);
 \item after multiplication by \((1-z)^{c/12}\), the same direct sewing
 converges through level five to the coefficients of \((1-z)^A\), with
 the component dependence vanishing as \(O(c^{-1})\).
\end{enumerate}

The checks are independent low-order certificates for the
global-family derivation.  In particular, they verify explicitly that
the \(G_0\)-ground-fibre and sign-dependent contributions combine into
the \(O(c^{-1})\) positive-quasi-primary remainder predicted by
Eq.~\eqref{eq:positive-quasiprimary-suppression}.

\end{document}
